%% kapa: exam=1 round=24 year=2013 period=2 form=A subject=accounting subject_ko=회계학 seq=14 number=14
%% source: Q-Net 감정평가사 기출문제 (cst00309) artlSeq=5073802 file=extracted/1008/5073802/2013년 제24회 감정평가사 1차시험 기출문제/2교시(A형).hwp
%% conversion: kordoc 4.24.51 -> tools/gen1.py (auto); review=proofread-sonnet; flags=none
\begin{question}{14}
수익인식에 관한 설명으로 옳지 않은 것은?
\begin{choices}
  \item 공장에서 이미 검사가 완료된 TV수상기의 설치와 같이 설치과정이 성격상 단순한 경우에는 구매자가 재화의 인도를 수락한 시점에 즉시 수익을 인식한다.
  \item 인도결제판매(cash on delivery sales)의 경우에는 인도가 완료되고 판매자가 현금을 수취할 때 수익을 인식한다.
  \item 프랜차이즈 수수료 중 계약에 의한 권리의 계속적인 사용에 부과되는 수수료나 계약기간동안 제공하는 기타 용역에 대한 수수료는 권리를 사용하는 시점이나 용역을 제공하는 시점에 수익으로 인식한다.
  \item 하나의 입장권으로 여러 행사에 참여할 수 있는 경우의 입장료수익은 각각의 행사를 위한 용역의 수행된 정도가 반영된 기준에 따라 각 행사에 배분하여 인식한다.
  \item 광고매체수수료는 광고 또는 상업방송이 대중에게 전달될 때 인식하고, 광고제작수수료는 제작이 완성되었을 때 인식한다.
\end{choices}
\end{question}
\answer{5}
